\section{Executive Summary}
\subsection{Overview}
A social engineering campaign was conducted against \reporttarget\ to measure susceptibility to
phishing and pretexting, and to identify departments/roles that would most benefit from targeted
awareness training.

\fillme{One or two sentences on the overall outcome, e.g. "18\% of recipients clicked the
simulated phishing link, and 6\% submitted credentials into the fake login page."}

\begin{criticalbox}
\textbf{Result --- \fillme{one-line verdict}.}\\[6pt]
Headline conclusions:
\begin{itemize}
  \item \fillme{Headline finding 1, e.g. "the Finance department had the highest click rate at 31\%"}
  \item \fillme{Headline finding 2, e.g. "only 4\% of recipients reported the email to IT/security"}
\end{itemize}
\end{criticalbox}

\subsection{Headline Metrics}
\begin{center}\begin{tabularx}{\linewidth}{@{}l r X@{}}
\toprule\rowcolor{TableHead}
\textcolor{white}{\bfseries Metric} & \textcolor{white}{\bfseries Value} & \textcolor{white}{\bfseries Benchmark} \\
\midrule
Emails delivered & \fillme{N} & --- \\
\rowcolor{TableAlt}
Open rate & \fillme{n\%} & \fillme{industry benchmark, if known} \\
Click rate & \fillme{n\%} & \fillme{industry benchmark} \\
\rowcolor{TableAlt}
Credential submission rate & \fillme{n\%} & \fillme{industry benchmark} \\
Report rate (to IT/security) & \fillme{n\%} & \fillme{industry benchmark} \\
\bottomrule\end{tabularx}\end{center}
